\documentclass{article}

\usepackage[utf8]{inputenc}
\usepackage{hyperref}
\usepackage[top=2cm]{geometry}

\usepackage[backend=biber, sorting=none]{biblatex}
\addbibresource{references.bib}

\title{%
  Analysing the Treewidth of Rust Programs \\
  \large Paper Proposal}

\author{
    Battini, Liam
    \and
    Haans, Geert
    \and
    Singh, Prabhjot
    \and
    Steffens, Tammo
}
\date{\today}

\begin{document}

\maketitle

\paragraph{Forming of the Idea}

Initially, we read a paper from our workgroup leader \cite{70be595fc13c49b985f1c3fe0868255b}, where we got the inspiration for our initial project idea.
The paper mentions constructs about a structured program and the bounded treewidth for goto-free languages and focuses on the control flow instead of the data flow, which inspired us to look for papers about treewidths and control flow.

\paragraph{Treewidth of Java Programs}

In the paper \textit{Treewidth of Java Programs}~\cite{gustedt_treewidth_2002}, there are two main contributions.
First, the authors present a subset of Java (label-free, which is similar to goto-free), which they relate to the analysis of the treewidth of goto-free C programs.
They show that when comparing the two subsets of Java and C, the results are the same.
Second, they present a tool to compute the treewidth of Java programs.
To show that in practice the upper bounds of 4-6 are rarely reached, they apply their tool to analyse the treewidth of selected parts of the Java standard library.

\paragraph{Control-Flow Graphs}
To analyse the structural properties of programs such as treewidth, we consider control-flow graphs (CFGs).
A CFG is a directed graph whose nodes correspond to basic blocks and whose edges represent possible transfers of control during execution~\cite{aho2006compilers}.
This representation captures the branching and looping structure of a program.
Accordingly, Gustedt et al~\cite{gustedt_treewidth_2002} analysed the treewidth of Java programs by computing the treewidth of their CFGs.
To study Rust in a comparable way, we therefore require an appropriate CFG representation of Rust programs.

\paragraph{Rust and MIR}
Rust is a structured programming language that does not support unrestricted \texttt{goto}-like control flow as described in the Rust language specifications~\cite{rustlang_loop_exprs}.
Given that it does not support \texttt{goto}-like control flow, it is feasible that we can make a clear comparison with the goto-free C language.
However, there may be other language constructs in Rust which could cause problems for bounding the treewidth.
We will investigate this.
To investigate whether Rust has other language constructs which may cause problems, we use the Mid-level Intermediate Representation (MIR) which is provided by the Rust compiler~\cite{rustc_mir} and is commonly used as a CFG-based representation of Rust programs in program verification tools such as Prusti~\cite{prusti}.
The MIR consists of basic blocks connected by explicit control-flow edges.
As such, MIR directly induces a CFG for each function.
This may allow us to extract the CFGs from Rust programs using an approach analogous to that of Gustedt et al~\cite{gustedt_treewidth_2002} and analyse their treewidth on a well-defined semantic representation.

\paragraph{Rust Control-Flow Graphs Analysis}
To analyse the treewidth of Rust programs, we need to extract the CFGs from Rust programs and compute their treewidth.
To extract the CFGs, we can use the MIR representation of Rust programs.
An existing tool, originally developed for data-flow-based vulnerability analysis of Rust programs ~\cite{Budde2024}, may be partially reusable for this purpose.

\paragraph{Goal}
For our project, we propose two possible contributions related to the analysis of the treewidth of Rust programs.
The first is a theoretical contribution that analyses the influence of the borrow checker on the treewidth of Rust programs and how to relate it to the analysis of goto-free C programs, similar to the analysis of the treewidth of label-free Java programs.
Selecting a suitable subset of Rust for this analysis is part of this contribution.
The second contribution is a practical one, where we implement a tool to compute the treewidth of Rust programs.
This tool will be used to analyse the treewidth in practice, e.g., analyse selected parts of the Rust standard library, to show that in practice the upper bounds of 4-6 are rarely reached.
Depending on our results and the time we have, the scope of the contributions may be adjusted to fit the project.

\vspace{1em}

\noindent\textit{Declaration: We used generative AI for reviewing and fixing grammar, finding related projects and making our writing more concise.}

\printbibliography

\end{document}
